% ============================================================================
% BACKUP of the Related Work section from main.tex, saved before shortening it.
% This is a verbatim dump for reference/rollback only; it is NOT \input by the
% paper. To restore, copy the block below back into main.tex over the current
% \section{Related Work} ... (just before \section{Method}).
% Saved: 2026-08-12.
% ============================================================================
%
% ----------------------------------------------------------------------------
% RUNNING CUT-LOG  (what was removed or de-emphasised while shortening RW)
% Full original text is preserved below this block, so anything here is
% recoverable. Reason codes: [intro]=belongs in intro (motivation),
% [method]=already covered in Method, [dup]=duplicated elsewhere in RW,
% [breadth]=tangential breadth cite, [detail]=fine-grained detail.
%
% PASS 1 (cuts + local merges; 8-paragraph skeleton kept)
%
% Citations fully CUT:
%   - madaan2023selfrefine (Self-Refine)      [breadth]  general self-refinement aside (P2)
%   - shinn2023reflexion   (Reflexion)        [breadth]  general self-refinement aside (P2)
%   - brown2024monkeys                        [breadth]  inference-scaling breadth (P3)
%   - snell2024scaling                        [breadth]  inference-scaling breadth (P3)
%   - geisler2025reinforce                    [dup]      redundant estimation cite (P6)
%   - arditi2024refusal   (abliteration)      [intro]    ** FLAG: tab:ablations still
%                                                          uses abliteration as a baseline;
%                                                          may want to re-cite it there **
%   - rimsky2024caa       (activation steering)[intro]   was only in the cut access sentence
%   - zou2023repe         (repr. engineering) [intro]    was only in the cut access sentence
%   - qi2024finetuning    (safety fine-tuning)[intro]    was only in the cut access sentence
%   (these four appeared ONLY in P5's "target-modifying constructions do not apply"
%    sentence and are now absent from the whole paper.)
%
% Citations DE-EMPHASISED (kept, but lost their own explanatory sentence;
% compressed into one clause):
%   - li2023contrastive (contrastive decoding)   (P4)
%   - liu2021dexperts   (DExperts)               (P4)
%   - sanchez2023cfg    (CFG-for-text)           (P4, now grouped with ho2022cfg)
%   - shi2024trusting   (context-aware decoding) (P4)
%
% Content CUT (no citation lost, prose removed):
%   - BLOOM four-stage enumeration                [method]  (P1 -> Method 3.x)
%   - "replace vanilla decoding with WILT" preview[method]  (P1 -> Method)
%   - FLRT logit/activation-difference-signal detail [detail] (P2)
%   - aranguri "10--300x more frequently" stat    [detail]  (P4)
%   - W2S "one set of weights in memory / co-loads"[detail] (P5)
%   - access-profile sentence: abliteration/activation-steering/
%     safety-finetuning "do not apply"            [intro]   (P5, made in intro/method)
%   - "same reason we don't benchmark white-box attacks" [dup] (P5, said in P2)
%   - greedy point-estimate "underestimates" lead [intro]   (P6, motivation)
%   - "how often vs. what it looks like" framing   [intro]  (P6 -> intro/discussion)
%   - WILT "never accesses weights..." restatement [intro]  (P8, made in intro/method)
%
% PASS 2 (structural fold/reorder: 8 paragraphs -> 4)
%
% New structure:
%   P1 "Automated behavioural auditing"  = old P1 (perez/BLOOM/Petri/multi-turn)
%                                          + old P8 (black-box debate) merged in
%   P2 "Input-side elicitation"          = old P2, unchanged
%   P3 "Output-side steering"            = old P3 (sampling/BoN) + old P4 (logit
%                                          arithmetic) + old P5 (constructing proposal)
%   P4 "Estimating vs. generating..."    = old P6 (estimation) + old P7 (chowdhury)
%
% W2S decision: the mechanistic contrast stays in Method ("Relation to W2S",
%   tied to eq:steer); RW P3 lists W2S as one construction and cross-refs Method.
%
% Citations: none newly CUT in Pass 2 (all preserved). aranguri2025diffamp was
%   cited TWICE pre-merge (old P4 + P5) -> now cited ONCE in P3's construction list.
%
% Prose CUT / de-emphasised in Pass 2 (recoverable from the full text below):
%   - old P8 "This favours our setting..." access restatement   [intro]  (setting is in intro)
%   - old P4 "CFG/context-aware apply it to prompts not models"  [detail] nuance dropped;
%       ho/sanchez/shi now listed only as product-of-experts family
%   - old P4 "We adapt the principle as an output-side mechanism"[dup]    (in intro + Method)
%   - old P5 "a single model...needing only inputs and outputs"  [dup]    (in intro + Method);
%       kept the shorter "no training, no second network"
%   - old P6 "...inspected and reused as adversarial-training data" [dup] now duplicated by the
%       intro sentence added this session; removed from RW
%   - "framework is agnostic to this choice" KEPT but shortened to one clause
% ----------------------------------------------------------------------------

\section{Related Work}\label{sec:related}

\paragraph{Automated behavioural evaluation.}
\citet{perez2022redteaming} introduced the LLM-evaluator-against-target paradigm, automating red-teaming by having one language model generate test inputs for another and a classifier score the resulting transcripts.
BLOOM~\citep{gupta2025bloom} formalises this into a four-stage pipeline (understanding, ideation, rollout, judgement) producing reproducible evaluation suites from a behaviour description.
Petri~\citep{fronsdal2025petri}, the most widely used auditing agent, instead takes an \emph{open-ended} approach (exploring for any concerning behaviour, with tools to set system prompts, roll back conversations, and prefill responses) rather than targeting a single named behaviour as we do.
\citet{huang2025multiturn} benchmark static, offline, and online methods for multi-turn behaviour elicitation, finding that static test cases saturate within a year while online methods keep finding failures.
We build on BLOOM: its \emph{multi-turn} rollouts probe behaviours in realistic conversations rather than single-turn prompts, the regime where many behavioural failures actually emerge~\citep{russinovich2024crescendo,anil2024manyshot}.
We retain its scaffold but replace the rollout's vanilla decoding with WILT, which steers the target's own next-token distribution toward the behaviour.


\paragraph{Input-side optimisation.}
A large body of work elicits target behaviours by optimising the \emph{input}.
Token-level attacks search for adversarial strings: GCG~\citep{zou2023universal} by greedy coordinate gradients, BEAST~\citep{sadasivan2024fast} by gradient-free beam search over the target's log-probabilities, and AutoDAN~\citep{liu2023autodan} by a genetic algorithm; FLRT~\citep{thompson2024flrt} optimises against a toxified fine-tune of the target, using its logit (or activation) difference from the safe model as the elicitation signal.
Prompt-iteration methods instead have an attacker model rewrite whole prompts from the target's replies (PAIR~\citep{chao2023jailbreaking}) or branch over them (TAP~\citep{mehrotra2023tap}); the same observe--revise--retry loop underlies self-refinement agents (Self-Refine~\citep{madaan2023selfrefine}, Reflexion~\citep{shinn2023reflexion}) and BLOOM's cross-round strategy refinement.
All of these spend their compute on the input side (searching tokens, iterating prompts, or refining strategies), and the gradient-based ones (GCG, FLRT) further assume the white-box access our setting lacks.
We find experimentally that, for behavioural elicitation, this is the weaker axis: steering the target's \emph{output} beats optimising its input, and layering input search or cross-round refinement on top adds little, so we keep cross-round refinement only as a baseline.


\paragraph{Sampling at attack time.}
Closely related to our output-side contribution, \citet{beyer2026sampling} reframe adversarial attacks as a resource-allocation problem between prompt optimisation and repeated output sampling, showing that integrating sampling boosts success rates by up to 37\% and efficiency by up to two orders of magnitude over greedy-decoding-only attacks.
\citet{hughes2024bon} push this further with Best-of-N jailbreaking, achieving high attack-success rates almost entirely through scaled output sampling over minimally augmented prompts.
These sit within a broader literature on inference-time compute scaling~\citep{brown2024monkeys,snell2024scaling} showing that repeated sampling reliably trades compute for task success across domains.
Their settings are single-turn attack benchmarks with unguided i.i.d.\ output sampling; we extend the idea to multi-turn behavioural rollouts and replace unguided sampling with a contrastive proposal, while sharing the core insight that input and output search consume the same compute budget and should be allocated jointly.


\paragraph{Output steering by logit arithmetic.}
Combining model distributions by adding their logits is a product of experts~\citep{hinton2002poe}, the principle behind a family of decoding-time steering methods.
Contrastive decoding~\citep{li2023contrastive} improves generation quality by subtracting amateur-model logits from expert-model logits; DExperts~\citep{liu2021dexperts} extends this for controllable generation with expert and anti-expert log-probability contributions.
Classifier-free guidance~\citep{ho2022cfg}, first developed for diffusion models and carried over to text~\citep{sanchez2023cfg}, and context-aware decoding~\citep{shi2024trusting} apply the same logit-interpolation principle to \emph{prompts} rather than models, amplifying the influence of an instruction or retrieved context at decoding time.
Closest to our use, \citet{aranguri2025diffamp} amplify the logit difference between post- and pre-training checkpoints ($\text{logits}_{\text{after}} + \alpha(\text{logits}_{\text{after}} - \text{logits}_{\text{before}})$) to surface rare misaligned behaviours $10$--$300\times$ more frequently.
We adapt the same principle as an \emph{output-side} red-teaming mechanism: a behaviour-conditioned copy of the target, obtained by system prompt rather than by training, serves as the second expert, steering sampling toward regions where undesirable continuations are concentrated without inflating the model's capability.


\paragraph{Constructing the unsafe proposal.}
Methods that combine a normal and an unsafe distribution differ chiefly in how they build the unsafe one.
\citet{angell2026tail} activation-steer it; weak-to-strong jailbreaking~\citep{zhao2024weak} and proxy-tuning~\citep{liu2024proxy} fine-tune a small proxy and transfer its safe-versus-unsafe logit difference to the target; \citet{aranguri2025diffamp} and emulated disalignment~\citep{zhou2024emulated} contrast a model against its pre-fine-tuning checkpoint; and \citet{chowdhury2025pathological} prompt the target with a jailbreak system prompt.
We take the prompting route, conditioning the unmodified target on a behaviour-eliciting system prompt: a single model, no training and no second network, requiring nothing but the ability to query the target and read its output token distribution, never its weights, activations, or gradients.
Because the second distribution comes from the target itself rather than a smaller proxy, it needs only \emph{one} set of weights in memory (unlike weak-to-strong, which co-loads the target with a separate safe--unsafe pair) and is as capable as the audited model: a small proxy may be too weak to exhibit a sophisticated behaviour at all. We test that claim directly in \cref{tab:ablations}, where the prompted term is swapped for each of these alternatives.
This access profile matters because auditing frontier models rarely permits more: their weights are not released, so target-modifying constructions (abliteration~\citep{arditi2024refusal}, activation steering~\citep{rimsky2024caa,zou2023repe}, or safety fine-tuning~\citep{qi2024finetuning}) do not apply.
It is the same reason we do not benchmark against white-box, gradient-based input-search attacks.
The framework is nonetheless agnostic to this choice: any of these constructions can supply the behaviour term when the access it assumes is available.


\paragraph{Distribution-level evaluation and rare-output estimation.}
A growing line of work argues that greedy, point-estimate evaluation systematically \emph{under}-estimates how often a model misbehaves, because undesired outputs live in the tail of the model's output distribution.
\citet{scholten2024probabilistic} give a formal probabilistic evaluation framework and show that state-of-the-art unlearning methods leak information once outputs are sampled rather than greedily decoded; \citet{beyer2026sampling} and \citet{geisler2025reinforce} make the analogous point for adversarial attacks, where sampling-aware and distributional objectives outperform greedy, single-affirmative-target ones.
Estimating these tail probabilities efficiently is itself a research problem: \citet{wu2024rare} compare importance sampling and activation extrapolation for low-probability output estimation, both of which beat naive sampling, and \citet{jones2025forecasting} forecast per-query elicitation probability from small-scale samples.
WILT sits squarely in this line, but with a different objective: rather than \emph{estimating} a propensity, we use a cheap output-side importance-sampling proposal (the behaviour-conditioned distribution) to \emph{surface concrete behaviour-exhibiting transcripts} in the multi-turn behavioural-evaluation setting.
A probability estimate says \emph{how often} a behaviour occurs; a concrete transcript shows \emph{what it looks like}, so it can be inspected to diagnose the underlying alignment failure and reused as training data for adversarial training that patches it.
And whereas the importance sampling of \citet{wu2024rare} searches over \emph{inputs} that give rise to a rare output, our proposal biases the \emph{output} distribution directly.
This mirrors \citet{angell2026tail}, who activation-steer an unsafe proposal model to importance-sample rare harmful outputs; the difference is that they reweight samples for an unbiased probability \emph{estimate}, whereas we induce the unsafe distribution by prompting and keep the steered samples as behaviour-exhibiting \emph{examples}.


\paragraph{Eliciting natural examples of rare behaviours.}
Closest to our goal, \citet{chowdhury2025pathological} likewise \emph{generate} concrete natural-language examples of rare pathological behaviours rather than estimating their probability.
They train reinforcement-learning ``investigator'' agents to craft prompts that elicit a target behaviour, optimising a propensity bound that, like WILT, is built on the contrast between the normal target and a jailbroken proposal (the target prompted with a jailbreak system prompt), keeping elicited behaviours close to the target's own distribution.
The differences are mechanism and cost: their elicitation is \emph{learned per behaviour} by RL (a separate run per rubric, hundreds of steps, ${\sim}10^5$ target samples, with manually added constraints) and acts on the \emph{input} prompt, whereas WILT is training-free, zero-shot per behaviour, and steers the \emph{output} logits directly.
We therefore do not compare head-to-head: their code is unreleased and the per-behaviour RL cost is prohibitive to reproduce faithfully, so we instead position WILT as a cheap, general alternative that reaches the same end (natural transcripts exhibiting a named rare behaviour) without learning a bespoke proposal for each one.


\paragraph{The black-box auditing regime.}
Whether auditing should rely on model internals is contested: \citet{casper2024blackbox} argue that black-box access yields only lower bounds on worst-case behaviour while white-box access tightens them, but \citet{sheshadri2026auditbench}, benchmarking auditing agents on models with implanted hidden behaviours, find the opposite in practice: ``the agent performs best with black-box tools''.
This favours our setting: WILT never accesses weights, activations, or gradients, using only the target's inputs and outputs.
Yet, unlike a static benchmark, it applies real optimisation pressure on the \emph{output} side, steering the target's own next-token distribution toward the behaviour and resampling to keep the highest-eliciting outputs (the inputs themselves are not tuned to each target), and so stays effective even against models that hide behaviours under direct questioning~\citep{vanderweij2025sandbagging}.
